\documentclass{ximera}

\title{Polynomial Identification}
\author{Kenneth Berglund}

\begin{document}
\begin{abstract}
    Identifying polynomials and finding their degrees and leading coefficients.
\end{abstract}
\maketitle

Look over the following table classifying expressions.

\begin{center}
\begin{tabular}{|c|c|c|c|} 
 \hline
 Expression & Polynomial? & Degree & Leading Coefficient \\
 \hline\hline
 \(x\) & Yes & 1 & 1 \\ 
 \hline
 \(3x^2\) & Yes & 2 & 3 \\
 \hline
 \(7x + x^{1/2}\) & No & N/A & N/A \\
 \hline
 \
 (2x^3 - x^2 + x - 4\) & Yes & 3 & 2 \\
 \hline
 \(\pi x^2 - 9\) & Yes & 2 & \pi \\
 \hline
 \(\frac{2}{3}x^5 + \frac{9}{7}x & Yes & 5 & \frac{2}{3} \\
 \hline
 \(\frac{2x^5 + 9x}{3x^5 - 7x} & No & N/A & N/A \\
 \hline
 \(x^\pi - 2\) & No & N/A & N/A \\
 \hline
 \(2^x + 2^x + 1\) & No & N/A & N/A \\
 \hline
 \((2 - x)(x + 1)\) & Yes & 2 & -1 \\
 \hline
\end{tabular}
\end{center}

\begin{enumerate}
    \item Come up with a definition of a polynomial. Reviewing the
    definition of a quadratic may help you. Check with your group
    and the instructor to confirm.
    \item Identify all of the following which are polynomials 
    and give their degree and leading coefficient.
    \begin{enumerate}
        \item \(f\left(x\right)=2x^3-x+1\)
	    \item \(g\left(x\right)=2x^\pi-x+1\)
	    \item \(h\left(x\right)=7-x\)
	    \item \(f\left(x\right)=x\)
	    \item \(g\left(x\right)=\left(x-17\right)\left(x+8\right)\)
	    \item \(h\left(x\right)=\left(2x-3\right)^5\left(x-1\right)^2\)
	    \item \(f\left(x\right)= e x^5-3\)
	    \item \(g\left(x\right)=\frac{x-1}{x+1}\)
	    \item \(h\left(x\right)=7x^6-2x+x^6\)
	    \item \(f\left(x\right)=1+x+x^2\)
	    \item \(g\left(x\right)=-3\left(x-\frac{1}{2}\right)^2+4\)
	    \item \(h\left(x\right)=\left(x^2+1\right)\left(x-1\right)\left(x+1\right)^2\)
	    \item \(f\left(x\right)=\left(-3x+2\right)\left(5x+1\right)\left(2x-2\right)^2\) 
    \end{enumerate}
\end{enumerate}


\end{document}
